\section{A direct PC simulation of the proof DAG}\label{sec:clause-simulation}

Unless stated otherwise, \(h\ge1\) and there are finitely many old
variables, say \(v\). Represent a linear clause by its affine true-indicators:
\[
 C=\bigvee_i(g_{C,i}=1),\qquad
 Z(C)=\{b:g_{C,i}(b)=0\text{ for all }i\}.
\]
An equation \(a(b)=\alpha\) has true-indicator \(1+a(b)+\alpha\).
The set \(Z(C)\) is an affine subspace or empty.

Assign to each nonempty clause a fresh complete accuracy-\(h\) block,
and write its product as \(P_C\). Set \(P_\varnothing=1\), using no
block for the empty clause. \lean{claims/AffineClauseRegistry.lean} Fix the whole block family before following
the proof DAG. Clause values \(P_C=0\) will be derived; they are not
included among its axioms.

\begin{lemma}[General product telescoping]\label{lem:telescoping}
\lean{claims/MpTelescoping.lean}
For polynomials \(g_i\), \(r_{ui}\) with finitely many indices \(i\), and \(h\ge0\), put
\[
 P_v=\prod_{u=0}^{v-1}\left(1-\sum_i r_{ui}g_i\right),\qquad
 U_{h,i}=\sum_{v=0}^{h-1}r_{vi}P_v.
\]
Then \(1-P_h=\sum_i U_{h,i}g_i\).
If \(\deg g_i\le\delta\), \(\deg r_{vi}\le1\), and \(h\ge1\), then
\[
 \deg P_h\le h(\delta+1),\qquad
 \deg U_{h,i}\le1+(h-1)(\delta+1),\qquad
 \deg(g_iP_h)\le\deg g_i+h(\delta+1).
\]
These are upper bounds; no freshness or nonvanishing of the inputs is assumed.
\end{lemma}
\begin{proof}
The identity holds at \(h=0\). Writing \(s_h=\sum_i r_{hi}g_i\), the
recurrences \(P_{h+1}=P_h(1-s_h)\) and
\(U_{h+1,i}=U_{h,i}+r_{hi}P_h\) give
\[
 1-P_{h+1}=1-P_h+P_hs_h=\sum_i U_{h+1,i}g_i.
\]
Each factor has degree at most \(\delta+1\), and each summand of
\(U_{h,i}\) has degree at most \(1+(h-1)(\delta+1)\).
The sum and product degree inequalities prove the three bounds. At \(h=0\), \(P_0=1\) and \(U_{0,i}=0\).
\end{proof}

\begin{corollary}[Affine-clause prefix identity]\label{lem:prefix}
\lean{claims/MpTelescoping.lean}
Every accuracy-\(h\) clause block, \(h\ge1\), has coefficients satisfying
\begin{equation}\label{eq:prefix}
 1-P_C=\sum_iU_{C,i}g_{C,i},\qquad
 \deg P_C\le2h,\qquad\deg U_{C,i}\le2h-1.
\end{equation}
Its companions have degree at most \(2h+1\).
\end{corollary}
\begin{proof}
Apply Lemma~\ref{lem:telescoping} with \(\delta=1\), reindexing coefficient
rows from zero-based to one-based. Exact companion degrees require the
additional hypotheses in Lemma~\ref{lem:cleanup}; a loose input-degree
bound alone would not justify equality.
\end{proof}

\begin{lemma}[Complementary-parity resolution]\label{lem:clause-resolution}\lean{claims/AffineClauseResolution.lean}
Suppose \(A=(\bigvee_i(a_i=1))\vee(u=1)\) and
\(B=(\bigvee_j(b_j=1))\vee(1-u=1)\), with \(u\) affine.
Let \(T\) have the union of the context inputs \(a_i,b_j\).
If \(P_A,P_B\) have PC derivations through degree \(K\), then \(P_T\)
has a derivation in the same fixed block system through
\(\max\{K,4h+1\}\).
\end{lemma}
\begin{proof}
The context companions \(a_iP_T,b_jP_T\) are axioms of the conclusion
block; a repeated input may use the same axiom.
The prefix identities yield
\begin{align}
 R_A&=P_T-U_{A,u}uP_T
       =P_AP_T+\sum_iU_{A,i}(a_iP_T),\label{eq:resolution-A}\\
 R_B&=P_T-U_{B,1-u}(1-u)P_T
       =P_BP_T+\sum_jU_{B,j}(b_jP_T).\label{eq:resolution-B}
\end{align}
Lemma~\ref{lem:pc-product} derives these through \(\max\{K,4h\}\).
In particular it multiplies the completed lines \(P_A,P_B\), not their
entire earlier derivations.

For an affine \(u=c+\sum_t c_tb_t\), the identity
\[
 u^2-u=\sum_t c_t(b_t^2-b_t)
\]
has a degree-two proof. Therefore
\begin{equation}\label{eq:pivot-correction}
 (1-u)R_A-U_{A,u}P_T(u^2-u)=(1-u)P_T
\end{equation}
is derivable through \(\max\{K,4h+1\}\).
The right side is a completed line of degree at most \(2h+1\).
Multiply it by \(U_{B,1-u}\), costing at most \(4h\), and add
\eqref{eq:resolution-B}. The result is \(P_T\).
All multiplications expand into primitive PC steps by
Lemma~\ref{lem:pc-product}.
\end{proof}

\begin{lemma}[Semantic weakening and tautologies]\label{lem:clause-weakening}\lean{claims/AffineClauseWeakening.lean}
If \(C\models D\) and \(P_C\) has a degree-\(K\) derivation, then
\(P_D\) is derivable through \(\max\{K,4h\}\).
A tautological \(D\) has a direct value derivation through \(2h+1\).
\end{lemma}
\begin{proof}
If \(Z(D)\) is nonempty, then \(Z(D)\subseteq Z(C)\).
Each \(g_{C,i}\) vanishes on \(Z(D)\), so affine linear algebra gives
constants \(\lambda_{ij}\) with
\[
 g_{C,i}=\sum_j\lambda_{ij}g_{D,j}.
\]
To see this directly, choose affine coordinates beginning with a basis
of the equations defining \(Z(D)\). An affine polynomial vanishing when
these coordinates are zero has only those coordinate terms.
Consequently \(g_{C,i}P_D\) is a constant linear combination of
conclusion companions. Now use
\[
 P_D=P_CP_D+\sum_iU_{C,i}(g_{C,i}P_D)
\]
and Lemma~\ref{lem:pc-product}, with ceiling \(\max\{K,4h\}\).

If \(Z(D)\) is empty, its affine equations are inconsistent, and their
span contains one. The same constant combination of their companions
gives \(P_D\), through \(2h+1\).
If the conclusion is empty, its value is one; the preceding formulas
still apply with no conclusion block. In the nonempty-zero-set case,
an implication to the empty clause forces all premise indicators to
be identically zero, so the already derived premise value is one.
\end{proof}


\begin{lemma}[Binary affine separator]\label{lem:binary-separator}
\lean{claims/BinaryAffineZeroCover.lean}
Let \(V\) be a vector space over \(\F\), and \(A,B,D\) finite clauses
of affine true-indicators. Suppose \(A\wedge B\models D\), but neither
\(A\models D\) nor \(B\models D\). There is an indicator \(u\) already
in \(A\), taking both values on \(W=Z(D)\), such that
\[
 W\cap Z(A)=\{x\in W:u(x)=0\},\qquad
 W\cap Z(B)=\{x\in W:u(x)=1\}.
\]
\end{lemma}
\begin{proof}
Choose \(a\in W\) satisfying \(A\) and \(b\in W\) satisfying \(B\).
Soundness forces \(a\in Z(B)\) and \(b\in Z(A)\). Choose indicators
\(u\in A\), \(v\in B\) with \(u(a)=v(b)=1\); then
\(u(b)=v(a)=0\). Soundness also gives \(W\subseteq Z(A)\cup Z(B)\).
If \(x\in W\cap Z(B)\) had \(u(x)=0\), the point \(a-b+x\) would
belong to \(W\) and satisfy both \(u=v=1\), a contradiction. Here we
used only the affine identity \(g(a-b+x)=g(a)-g(b)+g(x)\).
Thus \(u=1\) on \(W\cap Z(B)\), while \(u=0\) on \(W\cap Z(A)\).
The covering inclusion gives the two reverse implications and hence the
fiber equalities. The witnesses \(a,b\) give both values.
\end{proof}

\begin{lemma}[Two-premise semantic inference]\label{lem:affine-cover}
\lean{claims/BinarySemanticAffineCover.lean}
Every sound inference \(A\wedge B\models D\)
between finite affine clauses admits a derivation using at most three
semantic weakening or complementary-parity resolution steps, introducing
at most two auxiliary clauses. It is either a one-premise weakening or
has the form
\[
 A\models D\vee(u=1),\qquad B\models D\vee(u=0),
 \qquad\frac{D\vee(u=1)\quad D\vee(u=0)}D,
\]
where \(u\) can be chosen from \(A\). Semantic implication is preserved
under any affine restriction of the underlying space.
\end{lemma}
\begin{proof}
If either premise implies \(D\), use weakening. Otherwise apply
Lemma~\ref{lem:binary-separator}: on \(Z(D)\), satisfaction of \(A\)
forces \(u=1\), and satisfaction of \(B\) forces \(u=0\).
This proves the two weakenings; resolution concludes. Restriction simply
precomposes the clauses with an affine map and preserves every implication.
\end{proof}

\begin{lemma}[Clause compression]\label{lem:compression}
\lean{claims/AffineClauseCompression.lean}
On a finite-dimensional \(\F\)-vector space of dimension \(v\), every
finite affine clause has a semantically equivalent subclause of width at
most \(v+1\).
\end{lemma}
\begin{proof}
The space of affine forms has dimension \(v+1\). Extract a basis of the
span from the clause's own indicators. Simultaneous vanishing of the basis
is equivalent to vanishing of the entire span, hence gives the same
falsifying set. This works for tautologies as well as consistent systems
and introduces no new indicators.
\end{proof}

\paragraph{Idea of the registry theorem.}
The lemmas above simulate one inference at a time, given blocks for its
premises and its conclusion. A polynomial calculus refutation needs all of
its axioms fixed in advance, so the next theorem first reserves a block for
every clause that can occur: one for each node, two for the auxiliary
clauses that Lemma~\ref{lem:affine-cover} may introduce at that node, and
one for each supplied initial clause. It then walks through the DAG in
topological order. Because each step uses only the already derived lines
\(P_A\), \(P_B\) and costs at most \(4h+1\), the degree of the whole
derivation is the maximum, not the sum, of the local costs.

\begin{theorem}[A fixed registry for a finite proof DAG]\label{thm:dag-registry}
\lean{claims/AffineDAGRegistry.lean}
Consider an \(S\)-node affine-clause DAG on \(v\) old bits, using semantic
weakening, complementary resolution, or sound binary inference. Suppose
there is a finite supplied family \((A_j)_{j\in J}\), of widths at most
\(w\), such that each permitted initial clause is implied by some \(A_j\).
There is one complete accuracy-\(h\) registry with
\[
 N_{\mathrm{reg}}=3S+|J|\text{ slots},\qquad W=\max\{w,v+2\},
\]
at most \(N_{\mathrm{reg}}W\) inputs and at most \(v+hN_{\mathrm{reg}}W\) variables. If the supplied
initial values have PC derivations in this registry through
\(D\ge4h+1\), then every compressed source value has a derivation through
that same \(D\). An empty final clause yields a refutation.
\end{theorem}
\begin{proof}
Compress each source clause to width at most \(v+1\), preserving its
semantics, and choose the plans from Lemma~\ref{lem:affine-cover}.
Reserve one primary and two auxiliary slots per source node, plus one
slot per supplied initial. An auxiliary clause adds one pivot indicator
to a primary clause, so has width at most \(v+2\). Unused slots can
contain the empty clause and contribute no inputs or coefficients.
Each input has \(h\) distinct coefficient variables, giving the inventory.
Fix these slots and all their complete axioms first. In a topological
order, simulate each planned weakening or resolution with
Lemmas~\ref{lem:clause-resolution} and \ref{lem:clause-weakening}.
Initial coverage is another weakening. Reuse completed premise values;
the degree ceiling never acquires a factor depending on proof height.
\end{proof}

Fresh unconstrained propositional variables may be fixed to zero before
this construction. Semantic implication survives restriction; a restricted
constant pivot is a semantic weakening. This does not authorize extension
axioms defining additional predicates.

\begin{theorem}[Complete clause-to-PC transfer]\label{thm:clause-transfer}\lean{claims/BitPHPClauseTransfer.lean}
For arbitrary \(m,\ell\ge0\), a refutation of the usual bit-PHP CNF
on \(m\) rows and \(n=2^\ell\) labels with \(S\) nodes yields, for every
integer \(h\ge1\), a complete one-level old-affine family over
\(\mathcal Q_{m,\ell}\) with at most \(3S+\binom m2\) blocks and an ordinary-PC
refutation through
\begin{equation}\label{eq:simulation-degree}
 D=\max\{2h+\ell,4h+1\}.
\end{equation}
With \(v=m\ell\), \(N_{\mathrm{reg}}=3S+\binom m2\), and
\(W=\max\{\ell,v+2\}\), the total input count is at most \(N_{\mathrm{reg}}W\),
and the total variable count is at most \(v+hN_{\mathrm{reg}}W\). There is no proof-height restriction.
\end{theorem}
\begin{proof}
\lean{claims/BitPHPInitialBridge.lean} Introduce the compact initial linear clauses
\[
 C_{ii'}=\bigvee_{t=1}^\ell(b_{it}+b_{i't}=1).
\]
Each \(C_{ii'}\) semantically implies each label-specific initial clause
for the same pair in \eqref{eq:bit-php}: both rows encoding the same
label falsifies \(C_{ii'}\).
Use these \(\binom m2\) clauses as the supplied family in
Theorem~\ref{thm:dag-registry}; their widths are at most \(\ell\).

For a compact initial clause set \(g_t=b_{it}+b_{i't}\) and
\(T_t=\prod_{s<t}(1-g_s)\).
Its value has the ordinary identity
\begin{equation}\label{eq:initial-value}
 P_{C_{ii'}}=E_{ii'}P_{C_{ii'}}
                 +\sum_{t=1}^{\ell}T_t(g_tP_{C_{ii'}}).
\end{equation}
The first term uses the old degree-\(\ell\) collision axiom;
the other terms use the block's companions.
Every multiple has degree at most \(2h+\ell\).
This is an NS, hence PC, derivation at that degree.

Apply Theorem~\ref{thm:dag-registry} in the source proof's acyclic order.
Initial values are obtained from \eqref{eq:initial-value} and
Lemma~\ref{lem:clause-weakening}. Each remaining inference is simulated
by Lemmas~\ref{lem:clause-resolution} and \ref{lem:clause-weakening}.
Their premise values are actual earlier derived lines, so the same
global ceiling \eqref{eq:simulation-degree} persists along the DAG.
The final empty clause has value one.
All blocks have inputs affine in the old \(v\) bits and disjoint fresh
coefficients. Their full companion and coefficient-domain families
are present. No balancing or multiplication of an entire preceding
derivation is used.
\end{proof}
